% T24 -- Systemes embarques, reseaux de terrain et trames -- banque d'exercices
%
% ===================================================================================
% COMPLEMENTARITE AVEC LA FICHE fiches/T24_Systemes_Embarques.tex
% ===================================================================================
% Les trois Questions Types du Bloc 3 de la fiche sont :
%   QT1 = registre ADCON0 du PIC 18F8723, selection de la voie AN2 + GO/DONE
%         (= 2022, epreuve ecrite disciplinaire, question 25) ;
%   QT2 = PWM servomoteur 50 Hz, N = 1800 positions, f_ech = 90 kHz
%         (= 2025, epreuve ecrite disciplinaire appliquee, question 21) ;
%   QT3 = mise a l'echelle automate S7-1500, PIW8 / 27 648 -> [-20 ; +200] degC.
% Le Bloc 1 (§1.6) et le Bloc 2 de la fiche traitent en outre la TAILLE de la trame
% CAN (108 bits = 1+11+1+2+4+64+15+1+2+7), corrigee lors de l'audit 2026-08-24.
%
% AUCUN des six exercices ci-dessous ne refait ce travail intellectuel :
%   - pas de configuration de registre ADCON0 (ni d'aucun registre de CAN interne) ;
%   - pas de calcul de f_ech d'un PWM servomoteur ni de rapport cyclique ;
%   - pas de mise a l'echelle d'entree analogique d'automate ;
%   - pas de RECOMPTAGE de la taille d'une trame CAN a partir de sa structure
%     (c'est exactement la question D4 du sujet 2026, laissee a la fiche) ;
%     l'exercice 3 part au contraire des 111 bits DONNES par l'enonce du sujet et
%     porte sur la couche physique et le chronogramme, travail different.
%
% ===================================================================================
% SOURCES (toutes relues dans les .md de conversion Marker, chemin
%   docs/CAPET_Markdown-20260629T134311Z-3-001/CAPET_Markdown/docs/CAPET_SII_IE/
%   06_PDFs_Originaux/{annee}/{fichier}.md)
% ===================================================================================
%   - Ex. 1 : 2019, "EDE ENE 2" (code d'epreuve lu l.27 et l.656 du sujet ; le sujet
%     ne porte pas d'autre intitule exploitable, voir la note d'honnetete plus bas),
%     questions 15, 16 et 17. Contexte (reseau CANopen d'un bloc batteries de
%     locotracteur VMI) : E2_Sujet_2019.md l.322-324. Trame capturee
%     (ID 182h, DLC 05h, Data1..6 = 7Eh 01h F1h FFh BAh C2h) : l.324 (tableau).
%     Regles de decodage TPDO, formats Temp/Amp/Volt et les TROIS exemples de
%     verification (200,8 degC -> D8h/07h ; 100 A -> 64h/00h ; 515,34 V -> 4Eh/C9h) :
%     annexe A5, l.640-652. Structure de trame CAN 2.0A : l.612-636.
%     ATTENTION : le rapport de jury 2019 ne corrige PAS ces questions (sa numerotation
%     Q15-Q17 porte sur une autre partie, supercondensateurs, Rapport_Jury_2019.md
%     l.291-311). Les resultats ci-dessous sont donc MES calculs, valides par les
%     trois exemples de l'annexe A5 qui se retrouvent exactement.
%       -> noeud 182h - 180h = 2 ; T = 017Eh = 382 -> 38,2 degC ;
%          I = FFF1h (C2 sur 16 bits) = -15 -> -15 A donc CHARGE ;
%          U = C2BAh = 49 850 -> 498,50 V.
%
%   - Ex. 2 : 2023, "epreuve ecrite disciplinaire appliquee" (intitule lu l.13 du
%     sujet ; code EDE STI 2), questions 22, 23, 24 et 25.
%     Enonce (Modbus RTU sur RS 485, 102 ventilateurs / 51 par corridor, repeteurs) :
%     E2_Appliquee_Commune_2023.md l.343-371. Annexe 5 Modbus (maitre/esclave,
%     diffusion, codes fonction 03/06/10, CRC 16 bits calcule sur tous les octets
%     sauf lui-meme) : l.855-879 et l.893-905.
%     Corrige officiel Rapport_Jury_2023.md : Q22 l.740, Q23 l.744 (trame
%     00 06 0F A0 01 A2 22 C0), Q24 l.748 (01A2 = 418 -> 41,8 Hz), Q25 l.754 (ete).
%     Correspondance 41,8 Hz <-> 800 tr/min : Rapport_Jury_2023.md l.732 (Q20).
%     Citation jury : verbatim exact de Rapport_Jury_2023.md l.1039.
%     NB : la figure 11 (releve de la trame) est une image non extraite par Marker ;
%     la trame reproduite ici est celle du CORRIGE officiel, pas une reconstitution.
%
%   - Ex. 3 : 2026, "ADMISSIBILITE : EPREUVE 1" (intitule lu l.21 du sujet ; code
%     LDE ENE 1), partie D, questions D1, D2 et D5.
%     Enonce (bus CAN entre pile a combustible Symbio et calculateur Hycity) :
%     E1_Sujet_2026.md l.445-461. Couche physique differentielle CAN_H/CAN_L,
%     UD = UCANH - UCANL, seuils recessif/dominant, codage NRZ : l.471-481.
%     Exigence "message envoye en moins de 1 ms" et "trame constituee de 111 bits" :
%     l.513-517 (question D5). Le sujet donne donc 111 bits, PAS 108 : l'ecart de
%     3 bits est explicite dans le corrige (bits d'intermission ; voir remarque).
%     Le chronogramme DRD1 est une image ; son axe des temps extrait par Marker
%     porte la graduation 0-16-32-48-...-240 (E1_Sujet_2026.md l.752), soit un pas
%     de 16 us. Le chronogramme redessine ici respecte ce pas et conduit a
%     T_bit = 8 us / 125 kbit/s, seule valeur compatible a la fois avec ce pas et
%     avec l'exigence des 1 ms (111 x 8 = 888 us). Ce point est SIGNALE comme une
%     reconstruction dans le corrige, il n'est pas presente comme releve officiel.
%     PAS de rapport de jury 2026 disponible dans le corpus (dossier 2026 = sujets
%     E1 et E2 seulement) : aucune citation de jury n'est donnee sur cet exercice.
%
%   - Ex. 4 : 2024, "epreuve ecrite disciplinaire appliquee" (intitule lu l.13 du
%     sujet ; code EDE STI 2), questions 27, 28, 30, 31 et 32.
%     Enonce (Urbanloop, protocole UDP, codage hh:mm:ss, Tet = 20 us,
%     df = 300 Mbit/s) : E2_Appliquee_Commune_2024.md l.415-443.
%     Corrige officiel Rapport_Jury_2024.md l.771-792 : Q27 erreur 1 s -> 16,7 m ;
%     Q28 erreur > 1 cm, coder les ms ; Q30 40 octets = 320 bits, Ttrans = 1,07 us ;
%     Q31 Tret = 20 + 1,07 + 5 = 26,07 us ; Q32 erreur 0,435 mm, negligeable.
%     Le temps de decodage Td = 5 us n'apparait que dans le corrige officiel (Q31) :
%     il est donc donne explicitement dans l'enonce de l'exercice.
%
%   - Ex. 5 : 2023, "EPREUVE ECRITE DISCIPLINAIRE" (intitule lu l.5 du sujet, la
%     coquille "DSICIPLINAIRE" du PDF est corrigee ici ; option Ingenierie
%     electrique), questions E4, E7 et E8.
%     Enonce (robot viticole Bakus, module GPS, liaison a 115 200 bauds, vitesse
%     max 6 km/h, precision de position exigee 1 cm) : E2_Disciplinaire_IE_2023.md
%     l.486-520. Les figures 18 et 20 sont des images non extraites ; le NOMBRE de
%     caracteres (71) et le format (10 bits par caractere) proviennent du CORRIGE
%     officiel, Rapport_Jury_2023.md l.409, et sont donc donnes dans l'enonce de
%     l'exercice plutot que releves sur une figure reconstituee.
%     >> DEUX ERREURS DU CORRIGE OFFICIEL sont relevees et corrigees ici (E7/E8) :
%        (a) "710 bits soit 89 octets" : 710/8 = 88,75, et surtout les 710 bits
%            transmis ne sont pas 89 octets utiles -- il n'y a que 71 octets utiles
%            (568 bits), le reste etant de l'encapsulation ; c'est precisement
%            l'amalgame bit/octet que le jury reproche par ailleurs ;
%        (b) "la distance parcourue est 1,03 m" : la propre formule du rapport,
%            (6/3,6) x 710/115 200, vaut 0,0103 m = 1,03 CENTIMETRE. La conclusion
%            du rapport (cahier des charges non respecte) reste valide car
%            1,03 cm > 1 cm, mais de justesse et non d'un facteur 100.
%     Citation jury : verbatim exact de Rapport_Jury_2023.md l.463.
%
%   - Ex. 6 (docTech) : 2022, "EPREUVE ECRITE DISCIPLINAIRE" (intitule lu l.5 du
%     sujet, meme coquille "DSICIPLINAIRE" ; option Ingenierie electrique),
%     questions 16 et 17.
%     Enonce (poste de purge SNCF, TRT 200, capteur DS1820 sur bus OneWire) :
%     E1_Sujet_2022.md l.318-332 (= E2_Disciplinaire_IE_2022.md l.321-335, sujet
%     identique sur cette partie). Identifiant releve
%     (CC 3C 01 A8 16 48 C8 | 28) : figure 6, l.336. Sequence 0x44 / 750 ms / 0xBE :
%     l.352. Scratchpad releve (Byte 8..0 = E0 10 0C FF 7F 46 4B 03 36) : figure 7,
%     l.358. Extrait micropython "temp = Byte[1] << 8 | Byte[0] ; return temp / 16" :
%     figure 8, l.364-368. Table des operateurs micropython : DT C5, l.752-762.
%     Corrige officiel Rapport_Jury_2022.md l.246 (Q16 : 2^48 par famille, role du
%     CRC, famille DS18B20 car code family 28) et l.250-258 (Q17 : temp = 0x0336 =
%     822, concatenation, 12 bits car byte 4 = 7F, T = 51,375 degC, intervention
%     possible car < 60 degC).
%     >> Le contenu du docTech (DT C2 layout ROM 64 bits, DT C3 scratchpad 9 octets,
%     DT C4 format temperature et registre de configuration) correspond a des
%     figures du sujet NON extraites par Marker. Il est donc RECONSTITUE a partir
%     de la datasheet constructeur officielle DS18B20 (Dallas/Maxim, rev 042208),
%     verifiee page par page : ROM 64 bits = 8 bits CRC + 48 bits n de serie +
%     8 bits family code (28h) [fig. 4, p.6] ; scratchpad 9 octets, byte 0/1 =
%     temperature LSB/MSB, byte 4 = registre de configuration [p.8] ; registre de
%     configuration = 0 R1 R0 1 1 1 1 1, R1R0 = 11 -> 12 bits -> 0x7F [fig. 7, p.8] ;
%     temperature = mot signe de 16 bits en complement a 2, LSB = 2^-4 = 0,0625 degC
%     [table 2, p.5] ; commandes Convert T = 44h et Read Scratchpad = BEh [p.5].
%     Toutes ces valeurs coincident avec celles que le sujet et le corrige officiel
%     citent (byte 4 = 7F, division par 16, code family 28) : la reconstitution est
%     donc coherente avec le sujet, mais elle reste une reconstitution et est
%     annoncee comme telle dans le titre du docTech.
%     Source : https://cdn.sparkfun.com/datasheets/Sensors/Temp/DS18B20.pdf
%
% ===================================================================================
% NOTE D'HONNETETE sur les intitules d'epreuve
% ===================================================================================
% Le sujet 2019 ne contient, dans sa conversion Marker, aucun intitule d'epreuve en
% clair (les pages de garde sont des images) : seul le code administratif "EDE ENE 2"
% est lisible (l.27 et l.656). C'est donc ce code, et lui seul, qui est cite dans la
% reference de l'exercice 1 -- aucun intitule n'a ete suppose.
%
% ===================================================================================
% FPGA / VHDL
% ===================================================================================
% Recherche exhaustive (grep -ril "FPGA|VHDL|logique programmable") sur l'INTEGRALITE
% du corpus 2013-2026 : AUCUNE occurrence. Le constat du Bloc 4 de la fiche ("absent
% des sujets reels 2013-2026 mais au programme") est confirme. Aucun exercice de cette
% banque ne peut donc s'appuyer sur un sujet reel pour le VHDL : plutot que d'inventer
% une fausse reference, la notion n'est pas traitee ici et reste a reviser sur la fiche.
\expandafter\ifx\csname CAPETINCLUDEMODE\endcsname\relax
\documentclass[11pt,a4paper]{article}
% preamble.tex — packages et commandes communes a toutes les fiches CAPET SII IE
\input{glyphtounicode}
\pdfgentounicode=1
\usepackage{cmap}
\usepackage[T1]{fontenc}
\usepackage[utf8]{inputenc}
\usepackage[french]{babel}
% babel[french] active a \begin{document} des raccourcis de ponctuation (:;!?)
% qui cassent la syntaxe to[...] de circuitikz (erreur "+ or - expected") ;
% on les desactive, ce qui ne change que l'espacement fin avant ces signes.
\AtBeginDocument{\shorthandoff{:;!?}}
\usepackage[a4paper, top=1.6cm, bottom=1.6cm, left=1.4cm, right=1.4cm]{geometry}
\setlength{\headheight}{14.5pt}
\usepackage{amsmath, amssymb, amsfonts}
\usepackage{siunitx}
\sisetup{locale=FR, output-decimal-marker={,}, inter-unit-product=\ensuremath{{}\cdot{}},
         group-minimum-digits=4}
% Unites non fournies par siunitx, declarees ici une fois pour toutes (fiches + exercices).
% Les gardes \ifdefined evitent tout conflit avec une redeclaration locale restee dans un
% fichier : ne pas retirer les gardes, ne pas redeclarer ces unites dans les fichiers.
\DeclareSIUnit{\jour}{j}
\ifdefined\var\else\DeclareSIUnit{\var}{var}\fi   % puissance reactive
\ifdefined\vad\else\DeclareSIUnit{\vad}{VAD}\fi   % puissance deformante
\ifdefined\tr\else\DeclareSIUnit{\tr}{tr}\fi      % tour (tr/min)
\usepackage{array, booktabs, multirow, tabularx}
\usepackage{xcolor}
\usepackage[most, breakable]{tcolorbox}
\usepackage{multicol}
\usepackage{enumitem}
\usepackage{mathtools}
\usepackage{lmodern}
% Refonte visuelle 2026 (direction "moderne coloré") : Lato pour le texte,
% maths sans-serif assorties via newtxsf (charge apres amsmath/amssymb).
\usepackage[default]{lato}
\usepackage{newtxsf}
\usepackage{microtype}
\usepackage{fancyhdr}
\usepackage{tikz}
\usetikzlibrary{positioning,arrows.meta}
\usepackage{circuitikz}
\usepackage{lastpage}

% ── Couleurs par domaine (palette modernisée 2026) ────────────────────
\definecolor{colElec}{RGB}{29, 78, 216}        % Électrotechnique — bleu vif
\definecolor{colElecLight}{RGB}{219, 234, 254} % fond bleu clair
\definecolor{colAuto}{RGB}{21, 128, 61}        % Automatique — vert
\definecolor{colMeca}{RGB}{194, 65, 12}        % Mécanique — orange
\definecolor{colTherm}{RGB}{185, 28, 28}       % Thermique — rouge
\definecolor{colQT}{RGB}{232, 244, 233}        % (hérité, plus utilisé par questiontype)
\definecolor{colJury}{RGB}{255, 248, 225}      % fond jury (crème)
\definecolor{colJuryFrame}{RGB}{202, 138, 4}   % bordure/bandeau jury (ambre)
\definecolor{colGrayLight}{RGB}{245, 245, 247} % fond gris neutre
\definecolor{colSeparator}{RGB}{180, 180, 180} % gris séparateur
\definecolor{colExo}{RGB}{8, 145, 178}          % Banque d'exercices — cyan (distinct des fiches)
% ── Encre et gris de la refonte ───────────────────────────────────────
\definecolor{colInk}{RGB}{23, 32, 42}          % texte courant (quasi-noir)
\definecolor{colSlate}{RGB}{83, 98, 118}       % gris ardoise (méta, en-têtes)
\definecolor{colRule}{RGB}{203, 213, 225}      % filets clairs
\definecolor{colJuryBg}{RGB}{255, 248, 225}    % fond crème jury
\definecolor{colPresence}{RGB}{83, 98, 118}    % bloc 5 (présence dans les sujets)
\AtBeginDocument{\color{colInk}}

% ── Style encadré exercice + corrigé (banque d'exercices) ─────────────
\tcbset{
  exoBox/.style={
    breakable,
    enhanced,
    colback=colExo!8!white,
    colframe=colExo,
    colbacktitle=colExo,
    coltitle=white,
    fonttitle=\bfseries\small,
    left=4pt, right=4pt, top=3pt, bottom=3pt,
    before skip=5pt, after skip=5pt,
  }
}

% ── Bascule enonce / corrige (banque d'exercices) ─────────────────────
% \CAPETCORRIGE est defini en ligne de commande (pdflatex "\def\CAPETCORRIGE{1}\input{...}")
% pour compiler la variante corriges d'un meme fichier source ; par defaut (non defini),
% la compilation produit la variante enonces. \CAPETSHOW{enonce}{corrige} choisit la branche.
%
% TROIS MODES (refonte 2026-09-20, manuel de revision) :
%   - aucun drapeau   -> imprime l'enonce            (variante enonces, INCHANGE)
%   - \CAPETCORRIGE   -> imprime le corrige          (variante corriges, INCHANGE)
%   - \CAPETMANUEL    -> imprime l'enonce ET met le corrige de cote pour l'annexe
%
% Le mode manuel stocke chaque corrige dans une macro (\csgdef), sans ecriture
% fichier ni re-tokenisation : le contenu (circuitikz, accents, %) traverse intact.
% L'annexe les rejoue ensuite dans l'ordre via \CAPETANNEXE.
\usepackage{etoolbox}

\newcounter{CAPETexo}          % numero continu d'exercice dans le manuel
\setcounter{CAPETexo}{0}
\newcounter{CAPETannexe}       % compteur de relecture en annexe

% Titre courant de la notion en cours (pose par \CAPETNOTION), sert d'intertitre annexe
\newcommand{\CAPETnotioncode}{}
\newcommand{\CAPETnotiontitre}{}

% \CAPETNOTION{code}{titre} : ouvre une notion dans le manuel (memorise le libelle,
% et marque le rang du premier exercice de la notion pour l'intertitre d'annexe).
\newcommand{\CAPETNOTION}[2]{%
  \renewcommand{\CAPETnotioncode}{#1}%
  \renewcommand{\CAPETnotiontitre}{#2}%
  \csgdef{CAPETnotionstart\the\numexpr\value{CAPETexo}+1\relax}{#1 -- #2}%
}

% ── \degree robuste ───────────────────────────────────────────────────
% Certains fichiers d'exercices utilisent \degree. Ni le preambule ni siunitx
% ne le definissent, et \providecommand local ne suffit pas : sous hyperref,
% \degree existe deja comme commande factice (\ERROR), donc \providecommand
% s'abstient et l'appel echoue a la compilation du manuel (erreur invisible
% en compilation isolee). On le (re)definit donc inconditionnellement ici,
% apres le chargement des packages, pour tout le corpus.
\AtBeginDocument{\providecommand{\degree}{}\renewcommand{\degree}{\ensuremath{^\circ}}}

% Libelle decoratif : les fichiers d'exercices utilisent aussi \CAPETSHOW dans
% leur \fancyhead et leur bloc de titre, pour afficher "(enonces)" ou "(corriges)".
% Ces appels-la ne sont PAS des exercices : en mode manuel ils ne doivent ni
% consommer un numero, ni alimenter l'annexe. On les reconnait a leur premier
% argument, strictement egal au libelle ci-dessous.
\newcommand{\CAPETLIBELLEENONCES}{(\'enonc\'es)}

\ifdefined\CAPETMANUEL
  % --- Mode manuel : enonce en place, corrige differe vers l'annexe ---
  \newcommand{\CAPETSHOW}[2]{%
    \def\CAPETargun{#1}%
    \ifx\CAPETargun\CAPETLIBELLEENONCES
      #1%            appel decoratif : rendu tel quel, hors numerotation
    \else
      \CAPETSHOWEXO{#1}{#2}%
    \fi
  }
  \newcommand{\CAPETSHOWEXO}[2]{%
    \stepcounter{CAPETexo}%
    \csgdef{CAPETcorrige\theCAPETexo}{#2}%
    \edef\CAPETtmpnum{\theCAPETexo}%
    #1%
    \par\vspace{2pt}%
    {\footnotesize\itshape\color{colExo}%
     \hyperref[capet:cor:\CAPETtmpnum]{Corrige de l'exercice \CAPETtmpnum\ $\rightarrow$ p.~\pageref{capet:cor:\CAPETtmpnum}}%
     \label{capet:eno:\CAPETtmpnum}}%
  }
  % \CAPETANNEXE : rejoue tous les corriges captures, dans l'ordre, avec intertitres
  \newcommand{\CAPETANNEXE}{%
    \setcounter{CAPETannexe}{0}%
    \loop\ifnum\value{CAPETannexe}<\value{CAPETexo}%
      \stepcounter{CAPETannexe}%
      \edef\CAPETcurnum{\theCAPETannexe}%
      \ifcsname CAPETnotionstart\CAPETcurnum\endcsname
        \par\medskip
        {\normalsize\bfseries\color{colExo}\csuse{CAPETnotionstart\CAPETcurnum}}%
        \par\nobreak\smallskip
      \fi
      \begin{tcolorbox}[exoBox, breakable,
        title={Corrige de l'exercice \CAPETcurnum\
               \hfill\normalfont\footnotesize
               \hyperref[capet:eno:\CAPETcurnum]{enonce p.~\pageref{capet:eno:\CAPETcurnum}}}]%
      \label{capet:cor:\CAPETcurnum}%
      \csuse{CAPETcorrige\CAPETcurnum}%
      \end{tcolorbox}%
    \repeat
  }
\else
  % --- Modes historiques, comportement strictement inchange ---
  \newcommand{\CAPETSHOW}[2]{\ifdefined\CAPETCORRIGE #2\else #1\fi}
  \newcommand{\CAPETANNEXE}{}
\fi

% ── Difficulté 1 à 3 (étoiles pleines, amssymb) ───────────────────────
\newcommand{\difficulte}[1]{%
  \ifcase#1\relax
  \or $\bigstar$\;niveau 1/3%
  \or $\bigstar\bigstar$\;niveau 2/3%
  \or $\bigstar\bigstar\bigstar$\;niveau 3/3%
  \fi
}

% ── Style encadré extrait de document technique (banque d'exercices) ──
\tcbset{
  docTech/.style={
    breakable,
    enhanced,
    colback=colGrayLight,
    colframe=colSeparator,
    fonttitle=\bfseries\small\itshape,
    coltitle=black,
    colbacktitle=colGrayLight,
    left=4pt, right=4pt, top=3pt, bottom=3pt,
    before skip=4pt, after skip=4pt,
  }
}

% ── Style tcolorbox de section principale ────────────────────────────
% Refonte 2026 : bandeau titre coloré arrondi + fond très pâle, sans cadre.
\tcbset{
  ficheSection/.style={
    breakable, enhanced jigsaw,
    frame hidden, boxrule=0pt, arc=1.8mm,
    colback=#1!4!white,
    colbacktitle=#1,
    coltitle=white,
    fonttitle=\sffamily\bfseries,
    left=7pt, right=7pt, top=5pt, bottom=5pt,
    toptitle=3pt, bottomtitle=3pt,
    before skip=10pt, after skip=8pt,
  }
}

% ── Style encadré question type ──────────────────────────────────────
% Carte blanche, bandeau sombre, liseré très clair.
\tcbset{
  questiontype/.style={
    breakable, enhanced jigsaw,
    colback=white,
    frame hidden, boxrule=0pt, arc=1.8mm,
    colbacktitle=colElec!78!black,
    coltitle=white,
    fonttitle=\sffamily\bfseries\small,
    left=7pt, right=7pt, top=4pt, bottom=4pt,
    toptitle=2.5pt, bottomtitle=2.5pt,
    before skip=8pt, after skip=8pt,
    borderline={0.6pt}{0pt}{colElec!12!white},
  }
}

% ── Style encadré jury ───────────────────────────────────────────────
% Bandeau ambre, fond crème, sans cadre.
\tcbset{
  jury/.style={
    breakable, enhanced jigsaw,
    colback=colJuryBg,
    frame hidden, boxrule=0pt, arc=1.8mm,
    colbacktitle=colJuryFrame,
    coltitle=white,
    fonttitle=\sffamily\bfseries,
    left=7pt, right=7pt, top=4pt, bottom=4pt,
    toptitle=3pt, bottomtitle=3pt,
    before skip=10pt, after skip=8pt,
  }
}

% ── En-tête / pied de page communs (refonte 2026) ─────────────────────
% #1 code fiche, #2 titre court, #3 couleur domaine
\newcommand{\ficheHeaders}[3]{%
  \fancyhead[L]{\small\sffamily{\bfseries\textcolor{#3}{#1}}\hspace{0.7em}\textcolor{colSlate}{#2}}%
  \fancyhead[R]{\small\sffamily\textcolor{colSlate}{CAPET SII IE\, --- \,Session 2026}}%
  \fancyfoot[C]{\small\sffamily\textcolor{colSlate}{\thepage\,/\,\pageref{LastPage}}}%
  \fancyhead[C]{}\fancyfoot[L]{}\fancyfoot[R]{}%
}

% ── Bloc en-tête de fiche (BLOC 0, refonte 2026) ──────────────────────
% #1 code  #2 titre  #3 couleur domaine  #4 domaine  #5 ligne présence  #6 sous-thèmes
\newcommand{\ficheEntete}[6]{%
  \begin{tcolorbox}[enhanced, breakable,
    frame hidden, boxrule=0pt, arc=2.5mm,
    interior style={left color=#3, right color=#3!75!black},
    left=9pt, right=9pt, top=7pt, bottom=7pt,
    before skip=0pt, after skip=6pt]
    {\color{white}\sffamily\bfseries\LARGE #1\quad #2}\\[4pt]
    {\color{#3!22!white}\sffamily\small \MakeUppercase{#4}\quad\textbar\quad #5}\\[1pt]
    {\color{#3!35!white}\sffamily\footnotesize Sous-thèmes : #6}
  \end{tcolorbox}%
}

% ── Mise en page ─────────────────────────────────────────────────────
\pagestyle{fancy}
\fancyhf{}
\renewcommand{\headrulewidth}{0.5pt}
\renewcommand{\footrulewidth}{0pt}
\AtBeginDocument{%
  \renewcommand{\headrule}{{\color{colRule}\hrule width\headwidth height 0.5pt}}%
}

\setlist[itemize]{leftmargin=1.2em, itemsep=1pt, topsep=2pt, parsep=0pt}
\setlist[enumerate]{leftmargin=1.2em, itemsep=3pt, topsep=2pt, parsep=0pt}

\setlength{\parskip}{2pt}
\setlength{\columnsep}{10pt}

\def\CAPETBEGINFILE{\begin{document}}
\def\CAPETENDFILE{\end{document}}
\else
\def\CAPETBEGINFILE{}
\def\CAPETENDFILE{}
\fi

% Le baud (symbole/s) n'est pas fourni par siunitx ; garde \ifdefined comme pour \var.
\ifdefined\baud\else\DeclareSIUnit{\baud}{Bd}\fi

\fancyhead[L]{\textbf{\textcolor{colExo}{T24}} - Exercices - Syst\`emes embarqu\'es}
\fancyhead[R]{\small\textit{Banque d'exercices \CAPETSHOW{(\'enonc\'es)}{(corrig\'es)}}}
\fancyfoot[C]{\small\thepage\,/\,\pageref{LastPage}}

\CAPETBEGINFILE

\begin{tcolorbox}[
  enhanced, breakable,
  colback=colExo,
  colframe=colExo,
  left=6pt, right=6pt, top=5pt, bottom=5pt,
  before skip=0pt, after skip=4pt
]
  {\color{white}\bfseries\large T24 \quad Syst\`emes embarqu\'es -- r\'eseaux de terrain, trames et protocoles}\\[2pt]
  {\color{white}\small Banque d'exercices -- 6 exercices \CAPETSHOW{(\'enonc\'es)}{(corrig\'es)}}
\end{tcolorbox}

% =====================================================================
% Exercice 1 -- 2019, EDE ENE 2, questions 15, 16 et 17
% =====================================================================
\begin{tcolorbox}[exoBox,
  title={Exercice 1 -- \difficulte{1} \hfill\normalfont\footnotesize 2019, EDE ENE 2, questions 15, 16 et 17}
]
\CAPETSHOW{
Les blocs batteries d'un v\'ehicule de maintenance ferroviaire communiquent avec l'unit\'e de contr\^ole sur un r\'eseau \textbf{CANopen}. Chaque bloc batterie est associ\'e \`a un n\oe ud num\'erot\'e et transmet cycliquement une trame de donn\'ees \textbf{TPDO} (\emph{transmit-process data object}) dont l'identificateur vaut $\mathrm{180_{16}} + n^\circ\text{du n\oe ud}$.

Un traceur a captur\'e la trame suivante sur le r\'eseau (toutes les valeurs sont en hexad\'ecimal) :

\begin{center}
\footnotesize
\begin{tabular}{@{}lcccccccc@{}}
\toprule
& ID & DLC & Data1 & Data2 & Data3 & Data4 & Data5 & Data6 \\
\midrule
Valeur & $\mathrm{182_{16}}$ & $\mathrm{05_{16}}$ & $\mathrm{7E_{16}}$ & $\mathrm{01_{16}}$ & $\mathrm{F1_{16}}$ & $\mathrm{FF_{16}}$ & $\mathrm{BA_{16}}$ & $\mathrm{C2_{16}}$ \\
\bottomrule
\end{tabular}
\end{center}

Le document annexe du sujet donne la structure du champ de donn\'ees d'une trame TPDO de bloc batterie :

\begin{center}
\footnotesize
\begin{tabular}{@{}lcccccc@{}}
\toprule
Champ & Data1 & Data2 & Data3 & Data4 & Data5 & Data6 \\
\midrule
Contenu & Temp\_L & Temp\_H & Amp\_L & Amp\_H & Volt\_L & Volt\_H \\
\bottomrule
\end{tabular}
\end{center}

Les trois grandeurs sont transmises sur \textbf{deux octets en compl\'ement \`a 2}, l'octet \_L \'etant l'octet de poids \textbf{faible} et l'octet \_H l'octet de poids \textbf{fort} :

\begin{itemize}[leftmargin=1.5em, itemsep=1pt]
  \item la \textbf{temp\'erature} est transmise en \textbf{dixi\`emes de degr\'e}. Exemple donn\'e par le constructeur : une temp\'erature de \SI{200,8}{\degreeCelsius} est transmise avec $\text{Temp\_L} = \mathrm{D8_{16}}$ et $\text{Temp\_H} = \mathrm{07_{16}}$ ;
  \item le \textbf{courant} est compt\'e \textbf{positif lors de la d\'echarge}. Exemple : pour un courant de d\'echarge de \SI{100}{\ampere}, $\text{Amp\_L} = \mathrm{64_{16}}$ et $\text{Amp\_H} = \mathrm{00_{16}}$ ;
  \item la \textbf{tension} est transmise en \textbf{dizaines de millivolts}. Exemple : une tension de \SI{515,34}{\volt} est transmise avec $\text{Volt\_L} = \mathrm{4E_{16}}$ et $\text{Volt\_H} = \mathrm{C9_{16}}$.
\end{itemize}

\begin{enumerate}[leftmargin=1.5em]
  \item D\'eterminer le num\'ero de la batterie ayant transmis cette trame.
  \item D\'eterminer la valeur de la temp\'erature mesur\'ee. V\'erifier au pr\'ealable la r\`egle de reconstitution sur l'exemple fourni par le constructeur.
  \item D\'eterminer la valeur de la tension mesur\'ee.
  \item D\'eterminer la valeur du courant mesur\'e, et pr\'eciser si le bloc batterie est en charge ou en d\'echarge \`a l'instant de la capture. Conclure sur la coh\'erence de l'ensemble des trois mesures.
\end{enumerate}

}{
\textbf{Corrig\'e.}

\begin{enumerate}[leftmargin=1.5em]
  \item L'identificateur d'une trame TPDO vaut $\mathrm{180_{16}} + n$, donc :
  \[ n = \mathrm{182_{16}} - \mathrm{180_{16}} = \mathrm{02_{16}} = 2 \]
  La trame a \'et\'e \'emise par le \textbf{bloc batterie n\textsuperscript{o}~2}.

  \emph{Remarque.} Le DLC vaut $\mathrm{05_{16}} = 5$, alors que six octets de donn\'ees sont list\'es : c'est une incoh\'erence de l'\'enonc\'e du sujet, qu'un candidat a tout int\'er\^et \`a signaler sur sa copie. Elle ne bloque pas le d\'ecodage, les six octets \'etant fournis.

  \item \textbf{V\'erification de la r\`egle sur l'exemple constructeur.} Avec $\text{Temp\_L} = \mathrm{D8_{16}}$ (poids faible) et $\text{Temp\_H} = \mathrm{07_{16}}$ (poids fort), le mot de 16 bits reconstitu\'e est $\mathrm{07D8_{16}}$, soit $7\times 16^2 + 13\times 16 + 8 = 2\,008$, et $2\,008/10 = \SI{200,8}{\degreeCelsius}$ : la r\`egle est bien \guillemotleft\ octet de poids fort \`a gauche, r\'esultat en dixi\`emes de degr\'e \guillemotright.

  \smallskip
  \textbf{Application \`a la trame captur\'ee.} $\text{Temp\_L} = \mathrm{7E_{16}}$, $\text{Temp\_H} = \mathrm{01_{16}}$ :
  \[ \mathrm{017E_{16}} = 1\times 16^2 + 7\times 16 + 14 = 382 \qquad\Longrightarrow\qquad T = \frac{382}{10} = \boldsymbol{\SI{38,2}{\degreeCelsius}} \]
  Le bit de poids fort du mot vaut $0$ : la valeur est bien positive, aucune conversion de compl\'ement \`a 2 n'est \`a faire.

  \item $\text{Volt\_L} = \mathrm{BA_{16}}$, $\text{Volt\_H} = \mathrm{C2_{16}}$, donc le mot vaut $\mathrm{C2BA_{16}}$ :
  \[ \mathrm{C2BA_{16}} = 12\times 16^3 + 2\times 16^2 + 11\times 16 + 10 = 49\,152 + 512 + 176 + 10 = 49\,850 \]
  La tension \'etant transmise en dizaines de millivolts, c'est-\`a-dire par pas de \SI{10}{\milli\volt} :
  \[ U = 49\,850 \times \SI{10}{\milli\volt} = \boldsymbol{\SI{498,50}{\volt}} \]
  \emph{V\'erification sur l'exemple constructeur} : $\mathrm{C94E_{16}} = 51\,534$, et $51\,534\times\SI{10}{\milli\volt} = \SI{515,34}{\volt}$ : conforme.

  \item $\text{Amp\_L} = \mathrm{F1_{16}}$, $\text{Amp\_H} = \mathrm{FF_{16}}$, donc le mot vaut $\mathrm{FFF1_{16}}$. \textbf{Ici le bit de poids fort vaut 1} : le nombre est \textbf{n\'egatif} et doit \^etre interpr\'et\'e en compl\'ement \`a 2 sur 16 bits :
  \[ \mathrm{FFF1_{16}} = 65\,521 \qquad\Longrightarrow\qquad I = 65\,521 - 2^{16} = 65\,521 - 65\,536 = \boldsymbol{\SI{-15}{\ampere}} \]
  Le courant \'etant compt\'e \textbf{positif en d\'echarge}, un courant n\'egatif de \SI{-15}{\ampere} signifie que le bloc batterie n\textsuperscript{o}~2 est en \textbf{charge}, avec un courant de charge de \SI{15}{\ampere}.

  \smallskip
  \textbf{Conclusion.} Les trois mesures sont mutuellement coh\'erentes et d\'ecrivent un \textbf{bloc batterie en cours de recharge} : courant entrant de \SI{15}{\ampere}, tension de \SI{498,5}{\volt} (valeur \'elev\'ee, typique d'une fin de charge pour un pack de traction), temp\'erature de \SI{38,2}{\degreeCelsius} mod\'er\'ee mais sup\'erieure \`a l'ambiante, ce qui est l'\'echauffement attendu d'un pack en charge. C'est exactement l'information dont l'algorithme d'\'equilibrage a besoin : il peut poursuivre la charge de ce bloc, la temp\'erature restant loin des seuils d'alarme d'une chimie Li-ion. \textbf{Le point de m\'ethode d\'eterminant est le signe} : lire $\mathrm{FFF1_{16}}$ comme $65\,521$ au lieu de $-15$ aurait fait conclure \`a une d\'echarge de \SI{65}{\kilo\ampere}, valeur physiquement absurde qui doit alerter imm\'ediatement.
\end{enumerate}
}
\end{tcolorbox}

% =====================================================================
% Exercice 2 -- 2023, epreuve ecrite disciplinaire appliquee, Q22-Q25
% =====================================================================
\begin{tcolorbox}[exoBox,
  title={Exercice 2 -- \difficulte{2} \hfill\normalfont\footnotesize 2023, \'epreuve \'ecrite disciplinaire appliqu\'ee, questions 22 \`a 25}
]
\CAPETSHOW{
Une serre semi-ferm\'ee est ventil\'ee par \textbf{102 ventilateurs} (51 par corridor) entra\^in\'es par des machines asynchrones aliment\'ees par des modulateurs d'\'energie. Les consignes de fr\'equence sont transmises par l'ordinateur de supervision via un r\'eseau Ethernet, puis par une \textbf{liaison s\'erie multipoints RS 485} sur laquelle chaque passerelle joue le r\^ole de \textbf{ma\^itre} et les modulateurs celui d'\textbf{esclaves}, selon le protocole \textbf{Modbus RTU}.

\begin{center}
\begin{tikzpicture}[>=Stealth, scale=0.85, every node/.style={scale=0.85}]
  \node[draw, minimum width=1.9cm, minimum height=0.8cm, align=center, font=\scriptsize] (sup) at (0,1.6) {Supervision\\(Ethernet)};
  \node[draw, minimum width=1.7cm, minimum height=0.8cm, align=center, font=\scriptsize] (pas) at (3.0,1.6) {Passerelle\\\textbf{ma\^itre}};
  \draw[->] (sup.east) -- (pas.west);
  % bus RS485
  \draw[very thick, colExo] (1.3,0.4) -- (10.4,0.4);
  \node[font=\scriptsize, colExo, above left] at (10.4,0.4) {bus RS 485};
  \draw (pas.south) -- (3.0,0.4); \fill (3.0,0.4) circle (1.6pt);
  \foreach \x/\lab in {5.0/1, 6.6/2, 9.6/51} {
    \node[draw, minimum width=1.25cm, minimum height=0.7cm, align=center, font=\scriptsize] (e\lab) at (\x,-0.7) {Modul.\\esclave \lab};
    \draw (\x,0.4) -- (\x,-0.35); \fill (\x,0.4) circle (1.6pt);
  }
  \node[font=\small] at (8.1,-0.7) {$\cdots$};
  \node[draw, dashed, fill=white, minimum width=1.3cm, minimum height=0.5cm, font=\scriptsize] (rep) at (8.1,1.15) {r\'ep\'eteur};
  \draw[dashed] (8.1,0.9) -- (8.1,0.4); \fill (8.1,0.4) circle (1.6pt);
\end{tikzpicture}
\end{center}

Le protocole Modbus peut \^etre impl\'ement\'e sur une liaison de type \textbf{RS 232}, \textbf{RS 422} ou \textbf{RS 485}. On retiendra de la documentation du sujet que la liaison RS 232 est une liaison \textbf{point \`a point} (un seul ma\^itre, un seul esclave) et de port\'ee limit\'ee, tandis que la RS 422 admet \textbf{10 r\'ecepteurs} et la RS 485 \textbf{32 r\'ecepteurs} sur un m\^eme bus.

Le protocole Modbus propose notamment les fonctions suivantes, chacune cod\'ee sur un octet :

\begin{center}
\footnotesize
\begin{tabular}{@{}cl@{}}
\toprule
Code fonction (hexad\'ecimal) & Nom de la fonction \\
\midrule
$\mathrm{03}$ & Lecture de $N$ mots de sortie cons\'ecutifs \\
$\mathrm{06}$ & \'Ecriture d'un mot de sortie \\
$\mathrm{10}$ & \'Ecriture de $N$ mots de sortie \\
\bottomrule
\end{tabular}
\end{center}

Le ma\^itre peut s'adresser \`a un esclave donn\'e et attendre sa r\'eponse, ou \textbf{diffuser} une requ\^ete \`a tous les esclaves du bus (\emph{broadcast}), ceux-ci ex\'ecutant alors l'ordre \textbf{sans \'emettre de r\'eponse} ; l'adresse $\mathrm{00_{16}}$ est r\'eserv\'ee \`a la diffusion. Le CRC, calcul\'e sur 16 bits, porte sur tous les octets de la trame \`a l'exception de lui-m\^eme. Une trame Modbus RTU est constitu\'ee, dans l'ordre, de l'\textbf{adresse de l'esclave} (1 octet), du \textbf{code fonction} (1 octet), des \textbf{donn\'ees} et du \textbf{CRC} (2 octets).

Le relev\'e effectu\'e sur la liaison RS 485 d'un corridor donne la trame \'emise par la passerelle :

\begin{center}
\large\ttfamily 00\quad 06\quad 0F\quad A0\quad 01\quad A2\quad 22\quad C0
\end{center}

\begin{enumerate}[leftmargin=1.5em]
  \item Justifier le choix d'une liaison RS 485 parmi les trois liaisons propos\'ees, ainsi que la n\'ecessit\'e des r\'ep\'eteurs.
  \item Analyser la trame relev\'ee en indiquant l'adresse de l'esclave, la fonction, la donn\'ee et le CRC.
  \item Le registre Modbus $4000$ ($\mathrm{0FA0_{16}}$) des modulateurs contient la consigne de fr\'equence $f$ en Hz, avec un \textbf{facteur d'\'echelle 10} (la valeur r\'eelle est la valeur du registre divis\'ee par 10). Interpr\'eter ce message.
  \item L'\'etude m\'ecanique men\'ee en amont du sujet a \'etabli que les ventilateurs tournent \`a \SI{450}{\tr\per\minute} en hiver (soit $f = \SI{23,1}{\hertz}$) et \`a \SI{800}{\tr\per\minute} en \'et\'e (soit $f = \SI{41,8}{\hertz}$), un tiers seulement des ventilateurs fonctionnant en hiver contre la totalit\'e en \'et\'e. En d\'eduire le mode de fonctionnement de la serre \`a l'instant de l'\'emission de ce message. Justifier.
\end{enumerate}

}{
\textbf{Corrig\'e.}

\begin{enumerate}[leftmargin=1.5em]
  \item \textbf{\'Elimination de la RS 232.} Elle est \textbf{point \`a point} : elle n'autorise qu'un seul ma\^itre et un seul esclave, ce qui est incompatible avec une architecture ma\^itre/esclaves multipoints comportant 51 modulateurs par corridor ; sa port\'ee est en outre insuffisante \`a l'\'echelle d'un corridor de serre.

  \textbf{Choix entre RS 422 et RS 485.} Les deux sont des liaisons diff\'erentielles multipoints, mais la RS 485 admet \textbf{32 r\'ecepteurs} contre \textbf{10} pour la RS 422 : elle est donc nettement plus proche du besoin. La RS 485 utilise de surcro\^it un nombre de fils r\'eduit (une seule paire en half-duplex), ce qui diminue le co\^ut et la complexit\'e du c\^ablage sur une installation de cette taille.

  \textbf{N\'ecessit\'e des r\'ep\'eteurs.} M\^eme la RS 485 plafonne \`a 32 r\'ecepteurs, alors qu'il faut en adresser \textbf{51 par corridor} : $51 > 32$. Les r\'ep\'eteurs sont donc indispensables pour \textbf{\'etendre la limite en nombre d'\'equipements} (et accessoirement en distance) du segment, en r\'eg\'en\'erant le signal pour former plusieurs segments successifs.

  \item D\'ecoupage de la trame selon la structure adresse / fonction / donn\'ees / CRC :
  \begin{center}
  \footnotesize
  \begin{tabular}{@{}llll@{}}
  \toprule
  Octet(s) & Champ & Valeur & Signification \\
  \midrule
  1 & Adresse esclave & \texttt{00} & \textbf{diffusion} (\emph{broadcast}) : tous les esclaves \\
  2 & Code fonction & \texttt{06} & \'ecriture d'\textbf{un} mot de sortie \\
  3--4 & Donn\'ee & \texttt{0F A0} & adresse du registre vis\'e : $\mathrm{0FA0_{16}} = 4\,000$ \\
  5--6 & Donn\'ee & \texttt{01 A2} & valeur \'ecrite dans ce registre \\
  7--8 & CRC & \texttt{22 C0} & contr\^ole de redondance cyclique sur 16 bits \\
  \bottomrule
  \end{tabular}
  \end{center}
  Deux points m\'eritent d'\^etre soulign\'es sur la copie. D'une part l'adresse \texttt{00} n'est pas \guillemotleft\ l'esclave num\'ero z\'ero \guillemotright\ mais l'adresse r\'eserv\'ee de \textbf{diffusion} : le message s'adresse simultan\'ement \`a \textbf{tous} les modulateurs, et \textbf{aucun ne r\'epondra}. D'autre part le CRC ne fait pas partie des donn\'ees : il est calcul\'e sur les six octets pr\'ec\'edents et sert uniquement au destinataire \`a v\'erifier l'int\'egrit\'e de ce qu'il a re\c{c}u.

  \item La valeur \'ecrite dans le registre $4\,000$ est $\mathrm{01A2_{16}}$ :
  \[ \mathrm{01A2_{16}} = 1\times 16^2 + 10\times 16 + 2 = 256 + 160 + 2 = 418 \]
  Le facteur d'\'echelle \'etant de 10 :
  \[ f = \frac{418}{10} = \boldsymbol{\SI{41,8}{\hertz}} \]
  Le message est donc un ordre de \textbf{diffusion g\'en\'erale} imposant \`a \textbf{tous} les modulateurs d'\'energie de la ligne une consigne de fr\'equence statorique de \SI{41,8}{\hertz}.

  \item Deux indices convergent :
  \begin{itemize}
    \item la consigne vaut pr\'ecis\'ement \SI{41,8}{\hertz}, valeur qui correspond \`a \SI{800}{\tr\per\minute}, c'est-\`a-dire la vitesse \textbf{maximale}, celle du r\'egime d'\textbf{\'et\'e} (l'hiver correspondrait \`a \SI{23,1}{\hertz}) ;
    \item le message est \'emis en \textbf{diffusion} et impose la m\^eme consigne \`a tous les modulateurs, ce qui n'a de sens que si \textbf{tous} les ventilateurs fonctionnent -- or en hiver seul un tiers d'entre eux est en service.
  \end{itemize}

  \smallskip
  \textbf{Conclusion.} La serre fonctionne en \textbf{mode \'et\'e}, au \textbf{taux maximal de renouvellement d'air} : les 102 ventilateurs tournent simultan\'ement \`a \SI{800}{\tr\per\minute}. L'analyse ne s'arr\^ete donc pas au d\'ecodage de la trame : c'est le \textbf{croisement} de la valeur d\'ecod\'ee avec les r\'esultats de l'\'etude a\'eraulique \emph{et} avec le mode d'adressage (diffusion plut\^ot qu'adressage individuel) qui permet de qualifier l'\'etat r\'eel de l'installation et de valider que l'architecture r\'eseau retenue remplit bien sa fonction de pilotage.

  \smallskip
  \textit{Point de vigilance jury (rapport de jury CAPET 2023, commentaires de la partie 3) : \guillemotleft\ Le d\'ecodage d'une trame MODBUS n'a pas pos\'e de souci majeur. N\'eanmoins, l'interpr\'etation a \'et\'e plus d\'elicate. \guillemotright\ La question 2 ci-dessus est la partie facile ; ce sont les questions 3 et 4, qui demandent de \emph{donner un sens physique} aux octets, qui font la diff\'erence.}
\end{enumerate}
}
\end{tcolorbox}

% =====================================================================
% Exercice 3 -- 2026, epreuve 1, questions D1, D2 et D5
% =====================================================================
\begin{tcolorbox}[exoBox,
  title={Exercice 3 -- \difficulte{2} \hfill\normalfont\footnotesize 2026, \'epreuve 1 (mod\'elisation d'un produit), questions D1, D2 et D5}
]
\CAPETSHOW{
La pile \`a combustible d'un bus \`a hydrog\`ene communique avec le calculateur du v\'ehicule sur une architecture de type \textbf{bus CAN}. Le bus CAN est un bus s\'erie de terrain \textbf{multiplex\'e}, con\c{c}u pour fonctionner en environnement s\'ev\`ere (v\'ehicule, atelier, milieu industriel), dont l'\'enonc\'e rappelle les particularit\'es : bus producteur/consommateur o\`u tous les participants ont les m\^emes droits ; les n\oe uds ne sont \textbf{pas adress\'es}, seuls les \textbf{messages} portent un identificateur, chaque n\oe ud scrutant en permanence le bus pour d\'ecider si le message le concerne.

Dans sa version \emph{CAN High Speed}, la couche physique comporte deux signaux \textbf{CAN\_H} et \textbf{CAN\_L} construits de fa\c{c}on sym\'etrique (mode \textbf{diff\'erentiel}). C'est la diff\'erence de tension $U_D = U_{CAN\_H} - U_{CAN\_L}$ qui est prise en compte pour l'interpr\'etation des bits, avec un codage \textbf{NRZ} :

\begin{itemize}[leftmargin=1.5em, itemsep=1pt]
  \item si $U_D \leq \SI{0,5}{\volt}$ : \'etat \textbf{r\'ecessif}, \'etat logique \guillemotleft\ 1 \guillemotright\ ;
  \item si $U_D \geq \SI{0,9}{\volt}$ : \'etat \textbf{dominant}, \'etat logique \guillemotleft\ 0 \guillemotright.
\end{itemize}

Lors de la mise au point, un test de \textbf{proc\'edure d'urgence} est men\'e : la pile envoie aux autres \'equipements un message signalant un d\'efaut grave (surchauffe). Pour que la proc\'edure soit valid\'ee, le \textbf{message doit \^etre envoy\'e en moins de \SI{1}{\milli\second}}. Dans ce cas, \textbf{8 octets de donn\'ees} sont transmis, et l'\'enonc\'e pr\'ecise que la trame compl\`ete est alors constitu\'ee de \textbf{111 bits}.

Le chronogramme ci-dessous reproduit le d\'ebut de la trame captur\'ee sur le bus lors du test (la graduation de l'axe des temps est celle du document r\'eponse du sujet, en microsecondes) :

\begin{center}
\begin{tikzpicture}[>=Stealth, font=\footnotesize, xscale=0.345, yscale=0.85]
  % axes
  \draw[->] (0,0) -- (31.2,0) node[right] {$t$ (\si{\micro\second})};
  \draw[->] (0,-0.3) -- (0,2.45) node[above right, inner sep=1pt] {$U_D$ (\si{\volt})};
  % graduations toutes les 16 us
  \foreach \t in {0,16,32,48,64,80,96,112,128,144,160,176,192,208,224,240} {
    \pgfmathsetmacro\xx{\t/8}
    \draw (\xx,0) -- (\xx,-0.14) node[below, font=\scriptsize] {\t};
  }
  % seuils
  \draw[dashed, colSlate] (0,0.45) -- (30,0.45) node[right, font=\scriptsize, colSlate] {\SI{0,5}{\volt}};
  \draw[dashed, colSlate] (0,1.75) -- (30,1.75) node[right, font=\scriptsize, colSlate] {\SI{0,9}{\volt}};
  % signal : 1 bit = 1 unite x (=8us). Sequence : 0 1 1 0 1 0 0 0 1 1 0 1 0 1 1 0 ...
  % dominant (0 logique) -> UD haut (2.1) ; recessif (1 logique) -> UD bas (0.15)
  \def\hi{2.1}\def\lo{0.15}
  \draw[very thick, colExo]
    (0,\hi) -- (1,\hi) -- (1,\lo) -- (3,\lo) -- (3,\hi) -- (4,\hi) -- (4,\lo) --
    (5,\lo) -- (5,\hi) -- (8,\hi) -- (8,\lo) -- (10,\lo) -- (10,\hi) -- (11,\hi) --
    (11,\lo) -- (12,\lo) -- (12,\hi) -- (13,\hi) -- (13,\lo) -- (15,\lo) --
    (15,\hi) -- (16,\hi);
  \draw[colExo, dotted, very thick] (16,\hi) -- (19,\hi);
  \node[font=\scriptsize, colSlate, right] at (19.1,\hi) {suite de la trame};
  % reperage d'un bit : place SOUS l'axe des temps, a l'ecart du label U_D
  \draw[<->, colMeca, thick] (0,-0.62) -- (1,-0.62);
  \node[colMeca, font=\scriptsize, below] at (0.5,-0.66) {$T_{bit}$};
  \node[font=\scriptsize, left] at (-0.12,\hi) {dom.};
  \node[font=\scriptsize, left] at (-0.12,\lo) {r\'ec.};
\end{tikzpicture}
\end{center}

\begin{enumerate}[leftmargin=1.5em]
  \item Pour le protocole de communication CAN, pr\'eciser en justifiant : le \textbf{type} de transmission (s\'erie ou parall\`ele) ; la \textbf{nature} de la transmission (synchrone ou asynchrone) ; le \textbf{mode} de transmission (simplex, half-duplex ou full-duplex).
  \item Expliquer comment la couche physique diff\'erentielle permet l'\'elimination des perturbations \'electromagn\'etiques.
  \item \`A partir du chronogramme, d\'eterminer la \textbf{dur\'ee d'un bit} $T_{bit}$ ainsi que le \textbf{d\'ebit binaire} $D$ de la liaison.
  \item En d\'eduire la \textbf{dur\'ee de la trame} sachant qu'elle est constitu\'ee de 111 bits, puis conclure quant \`a l'exigence portant sur la dur\'ee de transmission.
\end{enumerate}

}{
\textbf{Corrig\'e.}

\begin{enumerate}[leftmargin=1.5em]
  \item \textbf{Type : transmission s\'erie.} Les bits sont transmis \textbf{les uns apr\`es les autres} sur une paire unique de conducteurs (CAN\_H / CAN\_L), et non simultan\'ement sur $n$ fils. C'est d'ailleurs ce qui rend n\'ecessaire le m\'ecanisme d'arbitrage bit \`a bit : un seul support est partag\'e.

  \textbf{Nature : transmission asynchrone.} Aucune horloge n'est transmise sur un conducteur d\'edi\'e. La resynchronisation des n\oe uds se fait \textbf{sur les transitions du signal lui-m\^eme}, notamment sur le front descendant du bit SOF en d\'ebut de trame. C'est pr\'ecis\'ement la raison d'\^etre du \emph{bit stuffing} du protocole CAN : l'insertion d'un bit de polarit\'e oppos\'ee apr\`es cinq bits identiques cons\'ecutifs garantit qu'il y aura toujours des transitions pour maintenir la synchronisation.

  \textbf{Mode : half-duplex.} La transmission est bien bidirectionnelle -- tout n\oe ud peut \'emettre comme recevoir -- mais \textbf{pas simultan\'ement} : le support est unique et partag\'e, et l'arbitrage garantit qu'un seul n\oe ud \'emet \`a la fois.

  \item Les deux conducteurs sont \textbf{physiquement torsad\'es et voisins} : une perturbation \'electromagn\'etique ext\'erieure induit sur chacun d'eux une tension parasite $u_p$ pratiquement \textbf{identique} (perturbation dite \emph{de mode commun}). Le r\'ecepteur n'exploitant que la \textbf{diff\'erence} :
  \[ U_D = (U_{CAN\_H} + u_p) - (U_{CAN\_L} + u_p) = U_{CAN\_H} - U_{CAN\_L} \]
  la perturbation \textbf{s'\'elimine par soustraction} et l'information binaire est restitu\'ee intacte. S'y ajoute la \textbf{marge de bruit} m\'enag\'ee par le constructeur : entre \SI{0,5}{\volt} et \SI{0,9}{\volt} s'\'etend une \textbf{zone morte} de \SI{0,4}{\volt} dans laquelle aucun \'etat n'est d\'eclar\'e, si bien qu'une perturbation r\'esiduelle de mode diff\'erentiel doit d\'epasser cette bande pour provoquer une erreur de lecture.

  \item Le chronogramme est gradu\'e tous les \SI{16}{\micro\second}, et l'on y compte \textbf{deux bits par graduation} : le palier \'el\'ementaire le plus court observ\'e -- par exemple le premier \'etat dominant, ou le bit r\'ecessif isol\'e -- a une largeur de \textbf{une demi-graduation}. D'o\`u :
  \[ T_{bit} = \frac{\SI{16}{\micro\second}}{2} = \boldsymbol{\SI{8}{\micro\second}} \]
  Le d\'ebit binaire est l'inverse de la dur\'ee d'un bit :
  \[ D = \frac{1}{T_{bit}} = \frac{1}{8\times 10^{-6}} = 125\,000\ \si{\bit\per\second} = \boldsymbol{\SI{125}{\kilo\bit\per\second}} \]
  C'est bien l'un des d\'ebits normalis\'es du bus CAN, et non une valeur quelconque, ce qui constitue une premi\`ere v\'erification de coh\'erence.

  \emph{M\'ethode.} Le pi\`ege classique consiste \`a relever la dur\'ee d'un \textbf{palier} au lieu de celle d'un \textbf{bit}. Ce sont deux choses diff\'erentes : le codage NRZ ne pr\'evoit \textbf{aucune transition entre deux bits identiques cons\'ecutifs}, si bien qu'un palier long correspond \`a plusieurs bits de m\^eme valeur. Il faut donc identifier le \textbf{plus petit} palier pr\'esent sur le relev\'e, qui seul vaut un bit unique.

  \item La trame comportant 111 bits :
  \[ T_{trame} = 111 \times T_{bit} = 111 \times 8 = \SI{888}{\micro\second} = \boldsymbol{\SI{0,888}{\milli\second}} \]

  \smallskip
  \textbf{Conclusion.} La dur\'ee de transmission du message d'urgence vaut \SI{0,888}{\milli\second}, \textbf{inf\'erieure aux \SI{1}{\milli\second}} impos\'ees par le cahier des charges : l'\textbf{exigence est satisfaite} et la proc\'edure d'urgence est valid\'ee sur ce crit\`ere. La marge est toutefois \textbf{faible -- \SI{11}{\percent} seulement}, et elle est de surcro\^it \textbf{optimiste} : le d\'ecompte de 111 bits ne tient pas compte des \textbf{bits de \emph{stuffing}} que le protocole ins\`ere apr\`es toute s\'equence de cinq bits identiques, ni d'une \'eventuelle attente d'arbitrage si un autre n\oe ud \'emet au m\^eme instant. Dans le pire cas, le \emph{stuffing} peut ajouter une vingtaine de bits, soit environ \SI{160}{\micro\second}, ce qui ferait franchir le seuil. La conclusion rigoureuse est donc : \textbf{exigence respect\'ee sur la trame nominale, mais sans marge suffisante pour absorber le \emph{stuffing}} -- il serait prudent de porter le d\'ebit \`a \SI{250}{\kilo\bit\per\second}, ce qui ram\`enerait la trame \`a \SI{444}{\micro\second} et offrirait un facteur 2 de s\'ecurit\'e.

  \smallskip
  \textbf{Remarque sur les 111 bits.} Le recomptage de la structure d'une trame CAN 2.0A standard \`a 8 octets de donn\'ees donne $1 + 11 + 1 + 2 + 4 + 64 + 15 + 1 + 2 + 7 = 108$ bits (SOF, identificateur, RTR, bits r\'eserv\'es, DLC, donn\'ees, CRC, d\'elimiteur, ACK, EOF). Les \textbf{3 bits} suppl\'ementaires annonc\'es par l'\'enonc\'e correspondent \`a l'\textbf{espace d'intertrame} (\emph{intermission}), qui s\'epare obligatoirement deux trames cons\'ecutives sur le bus. Les deux valeurs ne sont donc pas contradictoires : 108 bits pour la trame seule, 111 bits pour le temps d'occupation du bus -- c'est bien cette seconde valeur qu'il faut retenir pour comparer \`a un d\'elai impos\'e.
\end{enumerate}
}
\end{tcolorbox}

% =====================================================================
% Exercice 4 -- 2024, epreuve ecrite disciplinaire appliquee, Q27-Q32
% =====================================================================
\begin{tcolorbox}[exoBox,
  title={Exercice 4 -- \difficulte{2} \hfill\normalfont\footnotesize 2024, \'epreuve \'ecrite disciplinaire appliqu\'ee, questions 27, 28 et 30 \`a 32}
]
\CAPETSHOW{
Les capsules autonomes d'un syst\`eme de transport urbain de type \emph{Urbanloop} circulent \`a une vitesse maximale $V_0 = \SI{60}{\kilo\meter\per\hour}$. Chaque capsule \'elabore une information de position toutes les millisecondes et la transmet \`a la centrale de gestion en utilisant le protocole \textbf{UDP}, sur une \textbf{fibre optique} de d\'ebit $d_f = \SI{300}{\mega\bit\per\second}$.

Les premiers octets du champ \guillemotleft\ Data \guillemotright\ de la trame codent la \textbf{position} et l'\textbf{instant} associ\'e \`a cette position. Une premi\`ere proposition consiste \`a coder l'instant selon le syst\`eme horaire sur 24 heures \textbf{hh:mm:ss}, chacun des trois champs \'etant cod\'e sur un octet.

Le budget temporel de la cha\^ine de transmission comporte trois contributions :

\begin{itemize}[leftmargin=1.5em, itemsep=1pt]
  \item le temps d'\'elaboration de la trame, $T_{et} = \SI{20}{\micro\second}$ ;
  \item le temps de transmission $T_{trans}$ sur la fibre, \`a d\'eterminer ;
  \item le temps de d\'ecodage c\^ot\'e centrale, $T_d = \SI{5}{\micro\second}$.
\end{itemize}

La trame UDP compl\`ete envoy\'ee par une capsule comporte \textbf{40 octets}. Le cahier des charges impose une \textbf{pr\'ecision de position de \SI{1}{\centi\meter}}.

\begin{center}
\begin{tikzpicture}[>=Stealth, scale=0.9, every node/.style={scale=0.9}]
  \node[draw, minimum width=1.7cm, minimum height=0.85cm, align=center, font=\scriptsize] (cap) at (0,0) {Capsule\\\SI{60}{\kilo\meter\per\hour}};
  \node[draw, minimum width=1.8cm, minimum height=0.85cm, align=center, font=\scriptsize] (ela) at (2.9,0) {\'Elaboration\\$T_{et}{=}\SI{20}{\micro\second}$};
  \node[draw, minimum width=2.2cm, minimum height=0.85cm, align=center, font=\scriptsize] (fib) at (6.3,0) {Fibre optique\\$d_f{=}\SI{300}{\mega\bit\per\second}$};
  \node[draw, minimum width=1.8cm, minimum height=0.85cm, align=center, font=\scriptsize] (dec) at (9.6,0) {D\'ecodage\\$T_d{=}\SI{5}{\micro\second}$};
  \node[align=center, font=\scriptsize] (cen) at (12.0,0) {Centrale\\de gestion};
  \draw[->] (cap.east) -- (ela.west);
  \draw[->] (ela.east) -- (fib.west) node[midway,above,font=\scriptsize]{40 o};
  \draw[->] (fib.east) -- (dec.west);
  \draw[->] (dec.east) -- (cen.west);
  \draw[<->, colMeca, thick] (2.0,-0.75) -- (10.5,-0.75);
  \node[colMeca, font=\scriptsize, below] at (6.3,-0.75) {$T_{ret}$ : retard total entre la position r\'eelle et la position connue de la centrale};
\end{tikzpicture}
\end{center}

\begin{enumerate}[leftmargin=1.5em]
  \item Donner l'\textbf{erreur maximale}, en secondes, commise sur l'instant avec le codage hh:mm:ss. V\'erifier que l'erreur maximale ainsi faite sur l'estimation de la position de la capsule vaut \SI{16,7}{\meter}.
  \item Conclure sur le respect du cahier des charges avec ce codage. En cas de non-respect, proposer une solution de codage permettant d'obtenir une erreur de position conforme.
  \item Donner le nombre $N_b$ de bits contenus dans la trame, puis en d\'eduire le temps de transmission $T_{trans}$ sur la fibre.
  \item Calculer la dur\'ee $T_{ret}$ s\'eparant l'instant o\`u la capsule acquiert sa position de l'instant o\`u la centrale de gestion dispose de l'information.
  \item Conclure sur l'erreur de position finalement obtenue.
\end{enumerate}

}{
\textbf{Corrig\'e.}

\begin{enumerate}[leftmargin=1.5em]
  \item Le champ le plus fin du codage hh:mm:ss est la \textbf{seconde}. Deux instants s\'epar\'es de moins d'une seconde portent donc le \textbf{m\^eme} horodatage : l'erreur maximale sur l'instant vaut
  \[ \Delta t_{max} = \SI{1}{\second} \]
  Pendant cette seconde, la capsule, \`a sa vitesse maximale, parcourt :
  \[ V_0 = \SI{60}{\kilo\meter\per\hour} = \frac{60}{3{,}6} = \SI{16,7}{\meter\per\second} \qquad\Longrightarrow\qquad \Delta x_{max} = V_0 \times \Delta t_{max} = 16{,}7\times 1 = \boldsymbol{\SI{16,7}{\meter}} \]
  L'erreur de position annonc\'ee est bien retrouv\'ee.

  \item Le cahier des charges exige \SI{1}{\centi\meter}, soit \SI{0,01}{\meter}. Or $\SI{16,7}{\meter} \ggg \SI{0,01}{\meter}$ : le codage hh:mm:ss est \textbf{tr\`es largement insuffisant}, d'un facteur $16{,}7/0{,}01 \approx 1\,700$.

  Pour respecter l'exigence, il faut que la r\'esolution temporelle $\delta t$ v\'erifie $V_0\,\delta t \leq \SI{0,01}{\meter}$, soit :
  \[ \delta t \leq \frac{0{,}01}{16{,}7} \approx \SI{6e-4}{\second} = \SI{0,6}{\milli\second} \]
  Il faut donc \textbf{descendre au moins \`a la milliseconde}, et plus s\^urement aux \textbf{centaines de microsecondes}, en ajoutant au codage un champ de fraction de seconde. C'est coh\'erent avec le fonctionnement du syst\`eme, qui \'elabore de toute fa\c{c}on une position \emph{toutes les millisecondes} : horodater \`a la seconde revenait \`a jeter 999 informations sur 1000.

  \item La trame comporte 40 octets, soit :
  \[ N_b = 40\times 8 = \boldsymbol{320\ \text{bits}} \]
  \[ T_{trans} = \frac{N_b}{d_f} = \frac{320}{300\times 10^6} \approx \SI{1,07e-6}{\second} = \boldsymbol{\SI{1,07}{\micro\second}} \]

  \item Le retard total est la somme des trois contributions de la cha\^ine :
  \[ T_{ret} = T_{et} + T_{trans} + T_d = 20 + 1{,}07 + 5 = \boldsymbol{\SI{26,07}{\micro\second}} \]
  On notera que la transmission proprement dite ne repr\'esente que \SI{4}{\percent} du retard : ce sont l'\'elaboration de la trame et son d\'ecodage qui dominent largement le budget.

  \item Pendant ce retard, la capsule a continu\'e d'avancer \`a \SI{16,7}{\meter\per\second} :
  \[ \Delta x = V_0\,T_{ret} = \frac{60}{3{,}6}\times 26{,}07\times 10^{-6} \approx \SI{4,35e-4}{\meter} = \boldsymbol{\SI{0,435}{\milli\meter}} \]

  \smallskip
  \textbf{Conclusion.} L'erreur de position imputable \`a la cha\^ine de communication (\'elaboration, transmission, d\'ecodage) vaut \SI{0,435}{\milli\meter}, soit plus de \textbf{vingt fois moins} que le centim\`etre exig\'e par le cahier des charges : la solution de transmission est \textbf{tr\`es largement conforme} sur ce crit\`ere, et cette erreur est en pratique n\'egligeable devant les autres sources d'incertitude (glissement des roues, r\'esolution du codeur, positionnement absolu aux traverses). Le r\'esultat d\'eterminant de l'\'etude est donc ailleurs : c'est le \textbf{choix du codage de l'instant} (question 2) qui conditionnait le respect de l'exigence, pas la performance du support physique. Une fibre trente fois plus rapide n'aurait rien chang\'e \`a l'erreur de \SI{16,7}{\meter} du codage \`a la seconde -- illustration du fait qu'un budget d'erreur doit \^etre attaqu\'e par son terme dominant.
\end{enumerate}
}
\end{tcolorbox}

% =====================================================================
% Exercice 5 -- 2023, epreuve ecrite disciplinaire (IE), Q E4, E7, E8
% =====================================================================
\begin{tcolorbox}[exoBox,
  title={Exercice 5 -- \difficulte{3} \hfill\normalfont\footnotesize 2023, \'epreuve \'ecrite disciplinaire (option IE), questions E4, E7 et E8}
]
\CAPETSHOW{
Le robot viticole autonome \emph{Bakus} acquiert sa position au moyen d'un \textbf{module GPS} reli\'e \`a sa carte de contr\^ole par une \textbf{liaison s\'erie asynchrone} cadenc\'ee \`a \textbf{115\,200 bauds}. Les donn\'ees \'echang\'ees sont des \textbf{caract\`eres cod\'es en ASCII sur un octet}.

Le chronogramme ci-dessous reproduit la transmission d'\textbf{un} caract\`ere du module GPS vers la carte de contr\^ole. Le format de trame utilis\'e comporte, dans l'ordre : un \textbf{bit de start} \`a l'\'etat bas, les \textbf{8 bits de donn\'ees} transmis du poids faible au poids fort, et un \textbf{bit de stop} \`a l'\'etat haut.

\begin{center}
\begin{tikzpicture}[>=Stealth, font=\footnotesize, xscale=1.05, yscale=0.8]
  \def\hi{1.55}\def\lo{0.1}
  \draw[->] (-0.5,0) -- (11.4,0) node[right] {$t$};
  \draw[->] (-0.5,-0.25) -- (-0.5,2.2) node[above] {$V$};
  % repos haut, start bas, data 0 1 0 0 1 0 0 1 (LSB first) -> caractere 0b10010010 ? cf corrige
  % Sequence tracee : repos(hi) start(lo) puis bits LSB->MSB : 1 0 0 1 0 0 1 0 puis stop(hi)
  \draw[very thick, colExo]
    (-0.5,\hi) -- (0,\hi) -- (0,\lo) -- (1,\lo)          % start
    -- (1,\hi) -- (2,\hi)                                  % b0 = 1
    -- (2,\lo) -- (4,\lo)                                  % b1 = 0, b2 = 0
    -- (4,\hi) -- (5,\hi)                                  % b3 = 1
    -- (5,\lo) -- (7,\lo)                                  % b4 = 0, b5 = 0
    -- (7,\hi) -- (8,\hi)                                  % b6 = 1
    -- (8,\lo) -- (9,\lo)                                  % b7 = 0
    -- (9,\hi) -- (10.9,\hi);                              % stop + repos
  % separations verticales
  \foreach \x in {0,1,2,3,4,5,6,7,8,9,10} {
    \draw[dotted, colRule] (\x,-0.1) -- (\x,1.95);
  }
  % etiquettes des intervalles
  \node[font=\scriptsize, colMeca] at (0.5,2.0) {start};
  \foreach \i/\x in {0/1.5, 1/2.5, 2/3.5, 3/4.5, 4/5.5, 5/6.5, 6/7.5, 7/8.5} {
    \node[font=\scriptsize] at (\x,2.0) {$b_{\i}$};
  }
  \node[font=\scriptsize, colMeca] at (9.5,2.0) {stop};
  % valeurs logiques lues
  \foreach \v/\x in {1/1.5, 0/2.5, 0/3.5, 1/4.5, 0/5.5, 0/6.5, 1/7.5, 0/8.5} {
    \node[font=\scriptsize, colSlate] at (\x,-0.42) {\v};
  }
  \node[font=\scriptsize, left, colSlate] at (-0.55,\hi) {1};
  \node[font=\scriptsize, left, colSlate] at (-0.55,\lo) {0};
\end{tikzpicture}
\end{center}

L'ensemble des caract\`eres ASCII transmis par le module GPS forme une \textbf{trame compl\`ete de 71 caract\`eres}. Le robot se d\'eplace \`a une vitesse maximale de \SI{6}{\kilo\meter\per\hour}, et l'exigence de pr\'ecision sur la mesure de position du robot est fix\'ee \`a \textbf{\SI{1}{\centi\meter}}.

\begin{enumerate}[leftmargin=1.5em]
  \item \`A l'aide du chronogramme, d\'eterminer le \textbf{nombre de bits transmis} pour un caract\`ere ASCII. Identifier le caract\`ere transmis sur le chronogramme (on donne $\mathrm{49_{16}} = $ \guillemotleft\ I \guillemotright, $\mathrm{52_{16}} = $ \guillemotleft\ R \guillemotright, $\mathrm{92_{16}}$ hors table ASCII 7 bits).
  \item D\'eterminer le nombre total de \textbf{bits transmis} pour la trame compl\`ete, et le distinguer soigneusement du nombre d'\textbf{octets utiles} qu'elle v\'ehicule. En d\'eduire l'\textbf{efficacit\'e} de la transmission.
  \item En d\'eduire le \textbf{temps n\'ecessaire} pour transmettre cette trame compl\`ete.
  \item D\'eterminer la \textbf{distance maximale} parcourue par le robot pendant cet intervalle de temps. Conclure vis-\`a-vis de l'exigence de pr\'ecision de la mesure de position.
\end{enumerate}

}{
\textbf{Corrig\'e.}

\begin{enumerate}[leftmargin=1.5em]
  \item Le chronogramme montre, pour un seul caract\`ere : \textbf{1 bit de start} + \textbf{8 bits de donn\'ees} + \textbf{1 bit de stop}, soit
  \[ n_{bits/car} = 1 + 8 + 1 = \boldsymbol{10\ \text{bits}} \]
  pour \textbf{8 bits d'information utile} seulement.

  \textbf{Lecture du caract\`ere.} Les bits de donn\'ees sont transmis \textbf{du poids faible vers le poids fort}, et le chronogramme se lit $b_0 b_1 \ldots b_7 = 1\,0\,0\,1\,0\,0\,1\,0$. Il faut donc \textbf{inverser l'ordre} pour \'ecrire l'octet du MSB au LSB :
  \[ b_7 b_6 b_5 b_4 b_3 b_2 b_1 b_0 = 0100\;1001_2 = \mathrm{49_{16}} = 73_{10} \]
  soit le caract\`ere ASCII \guillemotleft\ \textbf{I} \guillemotright. \emph{Lire le chronogramme de gauche \`a droite sans inverser} aurait donn\'e $1001\,0010_2 = \mathrm{92_{16}}$, qui n'appartient m\^eme pas \`a la table ASCII 7 bits : c'est l'erreur la plus fr\'equente, et elle est ici auto-d\'etectable.

  \item \textbf{C'est la question o\`u se joue la distinction bit/octet.} La trame compl\`ete comporte 71 caract\`eres. Chaque caract\`ere occupant 10 bits \emph{sur la ligne} :
  \[ N_{bits\ transmis} = 71 \times 10 = \boldsymbol{710\ \text{bits}} \]
  En revanche, l'information \textbf{utile} v\'ehicul\'ee est de \textbf{71 octets}, soit :
  \[ N_{bits\ utiles} = 71 \times 8 = 568\ \text{bits} \]
  Les 710 bits transmis ne font donc \textbf{pas} 710/8 = 88,75 octets utiles : ils transportent \textbf{71 octets utiles} plus \textbf{142 bits d'encapsulation} (71 bits de start + 71 bits de stop). L'efficacit\'e de la transmission vaut :
  \[ Ef = \frac{N_{bits\ utiles}}{N_{bits\ transmis}} = \frac{568}{710} = \frac{8}{10} = \boldsymbol{0{,}80} \quad\text{soit}\quad \SI{80}{\percent} \]
  Autrement dit, \textbf{un bit sur cinq transmis sur la ligne ne porte aucune information} : c'est le prix de l'asynchronisme, qui oblige \`a resynchroniser le r\'ecepteur \`a chaque caract\`ere.

  \item Le d\'ebit \'etant de 115\,200 bauds, soit ici 115\,200 bits par seconde (un symbole = un bit sur cette liaison), et le temps de transmission se calculant sur les bits \textbf{r\'eellement transmis} :
  \[ T_{trame} = \frac{N_{bits\ transmis}}{D} = \frac{710}{115\,200} \approx \SI{6,16e-3}{\second} = \boldsymbol{\SI{6,16}{\milli\second}} \]
  \emph{Pi\`ege \`a \'eviter} : utiliser les 568 bits utiles aurait donn\'e \SI{4,93}{\milli\second}, soit une sous-estimation de \SI{20}{\percent} du temps d'occupation de la ligne. Le d\'ebit s'applique \`a \textbf{tout} ce qui transite, encapsulation comprise.

  \item La vitesse maximale du robot vaut :
  \[ v = \SI{6}{\kilo\meter\per\hour} = \frac{6}{3{,}6} \approx \SI{1,67}{\meter\per\second} \]
  Pendant la transmission de la trame, le robot parcourt donc :
  \[ d = v\,T_{trame} = \frac{6}{3{,}6}\times \frac{710}{115\,200} \approx \SI{1,03e-2}{\meter} = \boldsymbol{\SI{1,03}{\centi\meter}} \]

  \smallskip
  \textbf{Conclusion.} La position transmise par le module GPS est d\'ej\`a p\'erim\'ee de \SI{1,03}{\centi\meter} au moment o\`u la carte de contr\^ole ach\`eve de la recevoir, alors que le cahier des charges exige une pr\'ecision de \SI{1}{\centi\meter} : l'exigence \textbf{n'est pas respect\'ee}, m\^eme si elle ne l'est que de tr\`es peu (\SI{3}{\percent} de d\'epassement). Ce d\'epassement est enti\`erement imputable \`a la \textbf{dur\'ee de transmission}, ind\'ependamment de toute erreur du capteur lui-m\^eme. Deux leviers sont disponibles : \textbf{augmenter le d\'ebit} de la liaison (passer \`a \SI{230400}{\baud} ram\`enerait $d$ \`a \SI{0,51}{\centi\meter}, avec une marge confortable), ou \textbf{r\'eduire la taille de la trame} en ne transmettant que les champs NMEA effectivement exploit\'es par le robot. La premi\`ere solution est la plus imm\'ediate ; c'est celle que retient le corrig\'e officiel, qui conclut qu'\guillemotleft\ il faudra choisir une liaison plus rapide \guillemotright.

  \smallskip
  \textit{Note de lecture critique du corrig\'e officiel.} Le rapport de jury 2023 \'ecrit pour cette question : \guillemotleft\ La trame comporte 71 caract\`eres soit 710 bits (10 bits par caract\`ere) soit 89 octets \guillemotright, puis conclut \`a une distance de \guillemotleft\ 1,03 m \guillemotright. Deux inexactitudes s'y sont gliss\'ees. D'abord les \guillemotleft\ 89 octets \guillemotright\ ($710/8 = 88{,}75$) m\'elangent pr\'ecis\'ement bits transmis et octets utiles : la trame v\'ehicule \textbf{71} octets utiles, pas 89. Ensuite, la formule donn\'ee par le rapport lui-m\^eme, $\frac{6}{3{,}6}\times\frac{710}{115\,200}$, vaut \num[group-minimum-digits=5]{0,0103}\,\si{\meter}, soit \textbf{\SI{1,03}{\centi\meter} et non \SI{1,03}{\meter}}. La conclusion du rapport (cahier des charges non respect\'e) reste juste, mais elle l'est de justesse et non d'un facteur 100. \textbf{Savoir contr\^oler l'ordre de grandeur d'un r\'esultat, y compris face \`a un corrig\'e, est exactement la comp\'etence attendue} : un robot viticole qui d\'eriverait d'un m\`etre entre deux mesures GPS ne pourrait tout simplement pas suivre un rang de vigne.

  \smallskip
  \textit{Point de vigilance jury (rapport de jury CAPET 2023, commentaires de la cinqui\`eme partie) : \guillemotleft\ La cinqui\`eme partie permettait d'aborder l'\'electronique num\'erique et la communication par liaison s\'erie. Le jury constate beaucoup d'amalgames entre bit et octet et invite les candidats \`a revoir le principe d'une liaison s\'erie. \guillemotright}
\end{enumerate}
}
\end{tcolorbox}

% =====================================================================
% Exercice 6 -- docTech -- 2022, epreuve ecrite disciplinaire (IE), Q16-Q17
% =====================================================================
\begin{tcolorbox}[exoBox,
  title={Exercice 6 -- \difficulte{3} \hfill\normalfont\footnotesize 2022, \'epreuve \'ecrite disciplinaire (option IE), questions 16 et 17}
]
\CAPETSHOW{
La temp\'erature r\'egnant dans la chambre d'un \textbf{poste de purge} du r\'eseau ferroviaire doit imp\'erativement \^etre connue du technicien de maintenance avant toute intervention : un tuyau ayant c\'ed\'e y ferait r\'egner une pression et une temp\'erature dangereuses. Le seuil de s\'ecurit\'e impos\'e pour une intervention humaine est de \textbf{\SI{60}{\degreeCelsius}}.

Cette temp\'erature est mesur\'ee par un composant \textbf{DS1820} int\'egr\'e au module de t\'el\'etransmission TRT~200. Ce capteur dit \guillemotleft\ intelligent \guillemotright\ est reli\'e au microcontr\^oleur par un bus \textbf{OneWire} (bus Dallas, deux fils seulement : donn\'ees et masse). Le constructeur produit deux familles de ce composant :

\begin{itemize}[leftmargin=1.5em, itemsep=1pt]
  \item le \textbf{DS18B20}, valeur de temp\'erature cod\'ee et configurable entre 9 et 12 bits ;
  \item le \textbf{DS18S20}, valeur de temp\'erature cod\'ee sur 9 bits.
\end{itemize}

\begin{tcolorbox}[docTech, title={Document technique -- DS18B20 / DS18S20 sur bus OneWire (d'apr\`es la datasheet constructeur Dallas-Maxim, reconstitu\'ee)}]
\footnotesize
\textbf{1. Identifiant du composant -- ROM las\'eris\'ee de 64 bits.} Chaque composant porte un identifiant unique de 64 bits, organis\'e comme suit (poids forts \`a gauche) :

\begin{center}
\begin{tabular}{@{}ccc@{}}
\toprule
8 bits & 48 bits & 8 bits \\
\textbf{CRC} & \textbf{Num\'ero de s\'erie} & \textbf{Code famille} \\
\midrule
\multicolumn{3}{@{}l}{Code famille : $\mathrm{28_{16}}$ pour le DS18B20 ; $\mathrm{10_{16}}$ pour le DS18S20.} \\
\bottomrule
\end{tabular}
\end{center}

Le CRC est calcul\'e sur les 56 premiers bits et permet au ma\^itre de v\'erifier que l'identifiant a \'et\'e re\c{c}u sans erreur.

\smallskip
\textbf{2. M\'emoire \emph{Scratchpad} -- 9 octets.}

\begin{center}
\begin{tabular}{@{}cl@{}}
\toprule
Octet & Contenu \\
\midrule
Byte 0 & Temp\'erature, octet de poids \textbf{faible} (LSB) \\
Byte 1 & Temp\'erature, octet de poids \textbf{fort} (MSB) \\
Byte 2 & Seuil d'alarme haute $T_H$ \\
Byte 3 & Seuil d'alarme basse $T_L$ \\
\textbf{Byte 4} & \textbf{Registre de configuration} \\
Byte 5--7 & R\'eserv\'es (usage interne) \\
Byte 8 & CRC des huit octets pr\'ec\'edents \\
\bottomrule
\end{tabular}
\end{center}

\smallskip
\textbf{3. Registre de configuration (Byte 4).}

\begin{center}
\begin{tabular}{@{}cccccccc@{}}
\toprule
b7 & b6 & b5 & b4 & b3 & b2 & b1 & b0 \\
\midrule
0 & R1 & R0 & 1 & 1 & 1 & 1 & 1 \\
\bottomrule
\end{tabular}
\end{center}

\begin{center}
\begin{tabular}{@{}cccc@{}}
\toprule
R1 & R0 & R\'esolution & Quantum \\
\midrule
0 & 0 & 9 bits & \SI{0,5}{\degreeCelsius} \\
0 & 1 & 10 bits & \SI{0,25}{\degreeCelsius} \\
1 & 0 & 11 bits & \SI{0,125}{\degreeCelsius} \\
1 & 1 & \textbf{12 bits} & \textbf{\num[group-minimum-digits=5]{0,0625}\,\si{\degreeCelsius}} \\
\bottomrule
\end{tabular}
\end{center}

\smallskip
\textbf{4. Format de la temp\'erature.} La temp\'erature est stock\'ee dans les Byte 1:Byte 0 sous la forme d'un mot sign\'e de 16 bits en \textbf{compl\'ement \`a 2}, le bit de poids le plus faible valant $2^{-4}$ :

\begin{center}
\begin{tabular}{@{}cccccccc@{}}
\toprule
\multicolumn{8}{@{}l}{Byte 0 (LSB) :} \\
$2^{3}$ & $2^{2}$ & $2^{1}$ & $2^{0}$ & $2^{-1}$ & $2^{-2}$ & $2^{-3}$ & $2^{-4}$ \\
\midrule
\multicolumn{8}{@{}l}{Byte 1 (MSB) :} \\
S & S & S & S & S & $2^{6}$ & $2^{5}$ & $2^{4}$ \\
\bottomrule
\end{tabular}
\end{center}

$S$ : bits de signe ($S = 0$ temp\'erature positive, $S = 1$ temp\'erature n\'egative).

\smallskip
\textbf{5. Commandes utilis\'ees.} $\mathrm{44_{16}}$ : lancement d'une conversion de temp\'erature (dur\'ee jusqu'\`a \SI{750}{\milli\second} en r\'esolution 12 bits). $\mathrm{BE_{16}}$ : lecture du \emph{Scratchpad} (les 9 octets sont renvoy\'es).

\smallskip
\textbf{6. Op\'erateurs MicroPython.} \texttt{\&} et bit \`a bit ; \texttt{|} ou bit \`a bit ; \texttt{\textasciicircum} ou exclusif ; \texttt{\textasciitilde} non bit \`a bit ; \texttt{<{}<} d\'ecalage \`a gauche ; \texttt{>{}>} d\'ecalage \`a droite.
\end{tcolorbox}

Une lecture de l'\textbf{identifiant} du capteur install\'e dans le poste de purge \guillemotleft\ PstPurge102 \guillemotright\ donne :

\begin{center}
\footnotesize
\begin{tabular}{@{}lcccccccc@{}}
\toprule
Octet & \multicolumn{7}{c}{$\longleftarrow$ poids forts \dots\ poids faibles $\longrightarrow$} & \\
Valeur & $\mathrm{CC}$ & $\mathrm{3C}$ & $\mathrm{01}$ & $\mathrm{A8}$ & $\mathrm{16}$ & $\mathrm{48}$ & $\mathrm{C8}$ & $\mathrm{28}$ \\
\bottomrule
\end{tabular}
\end{center}

Le microcontr\^oleur envoie ensuite $\mathrm{44_{16}}$ au capteur, attend \SI{750}{\milli\second}, puis envoie $\mathrm{BE_{16}}$. Le contenu du \emph{Scratchpad} ainsi r\'ecup\'er\'e est :

\begin{center}
\footnotesize
\begin{tabular}{@{}lccccccccc@{}}
\toprule
N\textsuperscript{o} octet & Byte 8 & Byte 7 & Byte 6 & Byte 5 & Byte 4 & Byte 3 & Byte 2 & Byte 1 & Byte 0 \\
\midrule
Valeur & $\mathrm{E0}$ & $\mathrm{10}$ & $\mathrm{0C}$ & $\mathrm{FF}$ & $\mathrm{7F}$ & $\mathrm{46}$ & $\mathrm{4B}$ & $\mathrm{03}$ & $\mathrm{36}$ \\
\bottomrule
\end{tabular}
\end{center}

La fonction de lecture du capteur, \'ecrite en MicroPython, comporte les deux lignes suivantes :

\begin{center}
\begin{tcolorbox}[colback=colGrayLight, colframe=colSeparator, boxrule=0.4pt, left=6pt, right=6pt, top=4pt, bottom=4pt, width=0.72\linewidth]
\ttfamily\small
1:\quad temp = Byte[1] <{}< 8 | Byte[0]\\
2:\quad return temp / 16
\end{tcolorbox}
\end{center}

\begin{enumerate}[leftmargin=1.5em]
  \item \textbf{Relever} dans le document technique l'organisation des 64 bits de l'identifiant, puis en d\'eduire le nombre total de capteurs par famille pouvant th\'eoriquement \^etre fabriqu\'es avec un identifiant unique. \textbf{Donner} le r\^ole du CRC. \textbf{Identifier}, \`a partir de l'identifiant relev\'e, la famille \`a laquelle appartient ce composant.
  \item \textbf{Relever} dans le \emph{Scratchpad} l'octet qui fixe la r\'esolution de la conversion, et en d\'eduire sur combien de bits la temp\'erature est cod\'ee ici. Justifier la coh\'erence avec la temporisation de \SI{750}{\milli\second}.
  \item Donner la valeur de la variable \texttt{temp} de la ligne 1 \textbf{en binaire}, et nommer l'op\'eration informatique r\'ealis\'ee par cette ligne.
  \item D\'ecoder la temp\'erature r\'egnant dans le poste de purge, par \textbf{deux m\'ethodes} : en appliquant la ligne 2 du programme, puis en lisant directement les poids binaires du format de temp\'erature. Le technicien peut-il intervenir en toute s\'ecurit\'e ?
\end{enumerate}

}{
\textbf{Corrig\'e.}

\textbf{\'Etape 1 -- Extraction des donn\'ees du document technique.}
\begin{itemize}[leftmargin=1.5em, itemsep=1pt]
  \item Identifiant 64 bits $=$ 8 bits de CRC $+$ \textbf{48 bits de num\'ero de s\'erie} $+$ 8 bits de code famille ($\mathrm{28_{16}}$ pour le DS18B20, $\mathrm{10_{16}}$ pour le DS18S20).
  \item Le \emph{Scratchpad} comporte 9 octets, dont \textbf{Byte 0 et Byte 1} portent la temp\'erature (LSB puis MSB) et \textbf{Byte 4} est le \textbf{registre de configuration}.
  \item Registre de configuration : $0\;R1\;R0\;1\;1\;1\;1\;1$, avec $R1R0 = 11$ pour la r\'esolution 12 bits.
  \item Temp\'erature : mot sign\'e de 16 bits en compl\'ement \`a 2, LSB de poids $2^{-4}$.
\end{itemize}

\begin{enumerate}[leftmargin=1.5em]
  \item \textbf{Nombre de capteurs par famille.} Le num\'ero de s\'erie occupe \textbf{48 bits}, ind\'ependamment du code famille et du CRC qui ne sont pas des degr\'es de libert\'e. Le nombre de combinaisons possibles vaut donc :
  \[ N = 2^{48} = 281\,474\,976\,710\,656 \approx \boldsymbol{2{,}81\times 10^{14}} \text{ capteurs par famille} \]
  \textbf{R\^ole du CRC.} Le CRC (\emph{Cyclical Redundancy Check}) est un code de contr\^ole d'int\'egrit\'e : recalcul\'e par le ma\^itre sur les 56 bits re\c{c}us et compar\'e aux 8 bits transmis, il permet de \textbf{d\'etecter une alt\'eration de l'identifiant} lors de son transport sur le bus OneWire. Il ne corrige pas l'erreur : il signale qu'il faut relire.

  \textbf{Famille du composant.} Le code famille est l'octet de \textbf{poids faible}, c'est-\`a-dire le dernier de la lecture : $\mathrm{28_{16}}$. Le composant install\'e est donc un \textbf{DS18B20} (r\'esolution configurable de 9 \`a 12 bits), et non un DS18S20 (qui porterait $\mathrm{10_{16}}$).

  \emph{Attention au sens de lecture} : l'octet de t\^ete de la lecture, $\mathrm{CC_{16}}$, est le \textbf{CRC} et non le code famille. Le confondre conduirait \`a chercher une famille $\mathrm{CC_{16}}$ inexistante.

  \item L'octet qui fixe la r\'esolution est le \textbf{Byte 4}, registre de configuration, qui vaut ici $\mathrm{7F_{16}}$ :
  \[ \mathrm{7F_{16}} = 0111\;1111_2 \quad\Longrightarrow\quad b_7 = 0,\ \underbrace{R1 = 1,\ R0 = 1}_{b_6,\ b_5},\ b_4\ldots b_0 = 11111 \]
  $R1R0 = 11$ : la temp\'erature est donc cod\'ee sur \textbf{12 bits}, avec un quantum de $q = 2^{-4} = \num[group-minimum-digits=5]{0,0625}\,\si{\degreeCelsius}$.

  \textbf{Coh\'erence avec la temporisation.} Les \SI{750}{\milli\second} d'attente impos\'ees entre l'envoi de $\mathrm{44_{16}}$ et celui de $\mathrm{BE_{16}}$ correspondent pr\'ecis\'ement au \textbf{temps de conversion maximal en r\'esolution 12 bits}, la plus lente des quatre. Les deux informations se confirment mutuellement : le microcontr\^oleur attend bien le temps requis par la configuration effectivement charg\'ee dans le capteur.

  \item La ligne 1 reconstitue un mot de 16 bits \`a partir des deux octets du \emph{Scratchpad} :
  \[ \texttt{Byte[1] <{}< 8} \;=\; \mathrm{03_{16}} \ll 8 \;=\; \mathrm{0300_{16}} \qquad\text{puis}\qquad \mathrm{0300_{16}}\ \texttt{|}\ \mathrm{36_{16}} \;=\; \mathrm{0336_{16}} \]
  Soit, en binaire :
  \[ \texttt{temp} = \boldsymbol{0000\;0011\;0011\;0110_2} \qquad (= \mathrm{0336_{16}} = 822_{10}) \]
  L'op\'eration r\'ealis\'ee -- un d\'ecalage \`a gauche de l'octet de poids fort suivi d'un \guillemotleft\ ou \guillemotright\ bit \`a bit avec l'octet de poids faible -- est une \textbf{concat\'enation} des deux octets en un mot unique de 16 bits. Le \guillemotleft\ ou \guillemotright\ fonctionne ici comme une addition sans retenue, les 8 bits de poids faible de \texttt{Byte[1] <{}< 8} \'etant tous nuls.

  \item \textbf{M\'ethode 1 -- par la ligne 2 du programme.} Le bit de poids fort du mot vaut $0$ : la temp\'erature est \textbf{positive}, aucune conversion de compl\'ement \`a 2 n'est requise.
  \[ T = \frac{\texttt{temp}}{16} = \frac{822}{16} = \boldsymbol{\SI{51,375}{\degreeCelsius}} \]
  La division par 16 n'est pas arbitraire : elle traduit exactement le fait que le bit de poids le plus faible vaut $2^{-4} = 1/16$ de degr\'e.

  \textbf{M\'ethode 2 -- par lecture directe des poids.} En rempla\c{c}ant chaque bit \`a 1 par son poids dans le format du document technique, le mot $0000\;0011\;0011\;0110_2$ donne :
  \[ T = 2^{5} + 2^{4} + 2^{1} + 2^{0} + 2^{-2} + 2^{-3} = 32 + 16 + 2 + 1 + 0{,}25 + 0{,}125 = \SI{51,375}{\degreeCelsius} \]
  Les deux m\'ethodes concordent, ce qui valide \`a la fois la lecture du \emph{Scratchpad} et l'interpr\'etation du format.

  \smallskip
  \textbf{Conclusion.} La temp\'erature r\'egnant dans le poste de purge vaut \SI{51,375}{\degreeCelsius}, \textbf{inf\'erieure au seuil de \SI{60}{\degreeCelsius}} fix\'e pour une intervention humaine : le technicien \textbf{peut intervenir en s\'ecurit\'e}. La marge n'est toutefois que de \SI{8,6}{\degreeCelsius}, et rien n'indique dans cette mesure unique si la temp\'erature est stable ou en train de monter -- ce qui justifie pleinement l'architecture retenue par le constructeur : l'historique des mesures est enregistr\'e en base de donn\'ees et consult\'e par le technicien \textbf{avant} son d\'eplacement, de mani\`ere \`a juger d'une \textbf{tendance} et pas seulement d'un instantan\'e. On notera enfin que la r\'esolution de \num[group-minimum-digits=5]{0,0625}\,\si{\degreeCelsius} est ici tr\`es largement surdimensionn\'ee au regard d'une d\'ecision binaire prise sur un seuil de \SI{60}{\degreeCelsius} : une configuration 9 bits (\SI{0,5}{\degreeCelsius}) aurait suffi et aurait divis\'e par huit le temps de conversion, donc la consommation du module -- argument non n\'egligeable pour un \'equipement de t\'el\'etransmission aliment\'e de fa\c{c}on autonome.
\end{enumerate}
}
\end{tcolorbox}

\CAPETENDFILE
